% GENERATED FILE. Edit canonical lessons, then run npm run generate:books.
% canonical-source-hash: fnv1a64:38b0f836fb740e46
% canonical-lessons: BN-C424-sel, BN-C424-ophar, BN-C424-monda, BN-C424-chahida, BN-C424-othanama-kora

\chapter{Prices Up and Down}
\label{ch:bn-prices-up-and-down}
\hlchaptermodality{424}

\begin{chapteropening}
\textbf{By the end of this chapter:} \emph{I can spot a sale or an offer, and say demand and prices go up and down.}

Complete the last lesson of chapter 424: I can spot a sale or an offer, and say demand and prices go up and down.
\end{chapteropening}

\section[\texorpdfstring{\bn{সেল} (sel)}{সেল (sel)}]{\bn{সেল} (sel) — a sale (at reduced prices)}

\label{lesson:BN-C424-sel}



\begin{quote}
Before the new one: say the Bengali for a size (of clothes or shoes), then the Bengali for to palm off.
\end{quote}

\subsection*{You'll want to know: \bn{সেল}}
\textbf{\bn{সেল}} — \emph{sel} — \textquotedblleft{}a sale (at reduced prices)\textquotedblright{}.

\begin{quote}
\bn{পুজোর} \bn{আগে} \bn{দোকানে} \bn{সেল} \bn{চলে।}
\end{quote}

\begin{quote}
\emph{pujor āge dokāne sel chôle.}
\end{quote}

\begin{quote}
Before the Pujas the shops have sales on.
\end{quote}

\begin{cousinweb}
\textbf{\bn{সেল}} (\emph{sel}) is borrowed from English \emph{sale}, in the shop-window sense of prices cut for a time.
\end{cousinweb}

\begin{grammarlens}[title={What you've built}]
One more word for getting around.
\end{grammarlens}

\subsection*{Your turn}
\begin{itemize}
\raggedright
  \item \emph{Say it:} \emph{sel}
  \item \emph{Say it:} \emph{sel}, once more
  \item \emph{Say it:} \emph{pujor āge dokāne sel chôle.} — Before the Pujas the shops have sales on.
  \item \emph{Recall:} say the Bengali for a size (of clothes or shoes), then the Bengali for to palm off, then say \emph{sel} again
\end{itemize}

\begin{tcolorbox}[breakable,colback=teal!4,colframe=teal!35!black,title={Before you move on}]
What does \bn{সেল} mean? (\textquotedblleft{}A sale (at reduced prices)\textquotedblright{}.) Say it once more.
\end{tcolorbox}
\section[\texorpdfstring{\bn{অফার} (ôphār)}{অফার (ôphār)}]{\bn{অফার} (ôphār) — an offer; a special deal}

\label{lesson:BN-C424-ophar}



\begin{quote}
Before the new one: say the Bengali for to palm off, then the Bengali for a sale (at reduced prices).
\end{quote}

\subsection*{You'll want to know: \bn{অফার}}
\textbf{\bn{অফার}} — \emph{ôphār} — \textquotedblleft{}an offer; a special deal\textquotedblright{}.

\begin{quote}
\bn{দোকানে} \bn{এই} \bn{সপ্তাহে} \bn{চায়ের} \bn{ওপর} \bn{অফার} \bn{আছে।}
\end{quote}

\begin{quote}
\emph{dokāne ei sôptāhe chāyer opôr ôphār āchhe.}
\end{quote}

\begin{quote}
This week the shop has an offer on tea.
\end{quote}

\begin{cousinweb}
\textbf{\bn{অফার}} (\emph{ôphār}) is borrowed from English \emph{offer}.
\end{cousinweb}

\begin{grammarlens}[title={What you've built}]
One more word for getting around.
\end{grammarlens}

\subsection*{Your turn}
\begin{itemize}
\raggedright
  \item \emph{Say it:} \emph{ôphār}
  \item \emph{Say it:} \emph{ôphār}, once more
  \item \emph{Say it:} \emph{dokāne ei sôptāhe chāyer opôr ôphār āchhe.} — This week the shop has an offer on tea.
  \item \emph{Recall:} say the Bengali for to palm off, then the Bengali for a sale (at reduced prices), then say \emph{ôphār} again
\end{itemize}

\begin{tcolorbox}[breakable,colback=teal!4,colframe=teal!35!black,title={Before you move on}]
What does \bn{অফার} mean? (\textquotedblleft{}An offer; a special deal\textquotedblright{}.) Say it once more.
\end{tcolorbox}
\section[\texorpdfstring{\bn{মন্দা} (môndā)}{মন্দা (môndā)}]{\bn{মন্দা} (môndā) — a slump; slack trade}

\label{lesson:BN-C424-monda}



\begin{quote}
Before the new one: say the Bengali for a sale (at reduced prices), then the Bengali for an offer.
\end{quote}

\subsection*{You'll want to know: \bn{মন্দা}}
\textbf{\bn{মন্দা}} — \emph{môndā} — \textquotedblleft{}a slump; slack trade\textquotedblright{}.

\begin{quote}
\bn{গত} \bn{বছর} \bn{ব্যবসায়} \bn{মন্দা} \bn{ছিল।}
\end{quote}

\begin{quote}
\emph{gôtô bôchhôr bæbsāy môndā chhilô.}
\end{quote}

\begin{quote}
Last year business was slack.
\end{quote}

\begin{cousinweb}
\textbf{\bn{মন্দা}} (\emph{môndā}) goes back to Sanskrit \emph{manda}, \textquotedblleft{}slow, weak, slack\textquotedblright{} (Monier-Williams): trade that has gone slow.
\end{cousinweb}

\begin{grammarlens}[title={What you've built}]
One more word for getting around.
\end{grammarlens}

\subsection*{Your turn}
\begin{itemize}
\raggedright
  \item \emph{Say it:} \emph{môndā}
  \item \emph{Say it:} \emph{môndā}, once more
  \item \emph{Say it:} \emph{gôtô bôchhôr bæbsāy môndā chhilô.} — Last year business was slack.
  \item \emph{Recall:} say the Bengali for a sale (at reduced prices), then the Bengali for an offer, then say \emph{môndā} again
\end{itemize}

\begin{tcolorbox}[breakable,colback=teal!4,colframe=teal!35!black,title={Before you move on}]
What does \bn{মন্দা} mean? (\textquotedblleft{}A slump; slack trade\textquotedblright{}.) Say it once more.
\end{tcolorbox}
\section[\texorpdfstring{\bn{চাহিদা} (chāhidā)}{চাহিদা (chāhidā)}]{\bn{চাহিদা} (chāhidā) — demand (for goods)}

\label{lesson:BN-C424-chahida}



\begin{quote}
Before the new one: say the Bengali for an offer, then the Bengali for a slump.
\end{quote}

\subsection*{You'll want to know: \bn{চাহিদা}}
\textbf{\bn{চাহিদা}} — \emph{chāhidā} — \textquotedblleft{}demand (for goods)\textquotedblright{}.

\begin{quote}
\bn{শীতে} \bn{কম্বলের} \bn{চাহিদা} \bn{বেশি।}
\end{quote}

\begin{quote}
\emph{shīte kômbôler chāhidā beshi.}
\end{quote}

\begin{quote}
In winter there is more demand for blankets.
\end{quote}

\begin{cousinweb}
\textbf{\bn{চাহিদা}} (\emph{chāhidā}) is built on \bn{চাহ}-, the old stem of \textbf{\bn{চাওয়া}}, \textquotedblleft{}to want\textquotedblright{}, from the reconstructed \emph{cāh}, \textquotedblleft{}see, desire\textquotedblright{} (CDIAL 4775): what people want to buy.
\end{cousinweb}

\begin{grammarlens}[title={What you've built}]
One more word for getting around.
\end{grammarlens}

\subsection*{Your turn}
\begin{itemize}
\raggedright
  \item \emph{Say it:} \emph{chāhidā}
  \item \emph{Say it:} \emph{chāhidā}, once more
  \item \emph{Say it:} \emph{shīte kômbôler chāhidā beshi.} — In winter there is more demand for blankets.
  \item \emph{Recall:} say the Bengali for an offer, then the Bengali for a slump, then say \emph{chāhidā} again
\end{itemize}

\begin{tcolorbox}[breakable,colback=teal!4,colframe=teal!35!black,title={Before you move on}]
What does \bn{চাহিদা} mean? (\textquotedblleft{}Demand (for goods)\textquotedblright{}.) Say it once more.
\end{tcolorbox}
\section[\texorpdfstring{\bn{ওঠানামা} \bn{করা} (oṭhānāmā kôrā)}{ওঠানামা করা (oṭhānāmā kôrā)}]{\bn{ওঠানামা} \bn{করা} (oṭhānāmā kôrā) — to go up and down; to fluctuate}

\label{lesson:BN-C424-othanama-kora}



\begin{quote}
Before the new one: say the Bengali for a slump, then the Bengali for demand (for goods).
\end{quote}

\subsection*{You'll want to know: \bn{ওঠানামা} \bn{করা}}
\textbf{\bn{ওঠানামা} \bn{করা}} — \emph{oṭhānāmā kôrā} — \textquotedblleft{}to go up and down; to fluctuate\textquotedblright{}.

\begin{quote}
\bn{বাজারে} \bn{মাছের} \bn{দাম} \bn{রোজ} \bn{ওঠানামা} \bn{করে।}
\end{quote}

\begin{quote}
\emph{bājāre māchher dām roj oṭhānāmā kôre.}
\end{quote}

\begin{quote}
Fish prices in the market go up and down every day.
\end{quote}

Past (he/she …): \bn{সেই} \bn{দাম} \bn{ওঠানামা} \bn{করল}, \emph{sei dām oṭhānāmā kôrlô}.

\bn{ওঠানামা} alone is the rise and fall: \bn{দামের} \bn{ওঠানামা}, the ups and downs of prices; \bn{সিঁড়িতে} \bn{ওঠানামা}, going up and down stairs.

\begin{cousinweb}
\textbf{\bn{ওঠানামা} \bn{করা}} (\emph{oṭhānāmā kôrā}) pairs \textbf{\bn{ওঠা}}, \textquotedblleft{}to rise\textquotedblright{} (CDIAL 1900), with \textbf{\bn{নামা}}, \textquotedblleft{}to come down\textquotedblright{} (CDIAL 6964), and makes a verb of the pair with \bn{করা}.
\end{cousinweb}

\begin{grammarlens}[title={What you've built}]
One more word for getting around.
\end{grammarlens}

\subsection*{Your turn}
\begin{itemize}
\raggedright
  \item \emph{Say it:} \emph{oṭhānāmā kôrā}
  \item \emph{Say it:} \emph{oṭhānāmā kôrā}, once more
  \item \emph{Say it:} \emph{bājāre māchher dām roj oṭhānāmā kôre.} — Fish prices in the market go up and down every day.
  \item \emph{Recall:} say the Bengali for a slump, then the Bengali for demand (for goods), then say \emph{oṭhānāmā kôrā} again
\end{itemize}

\begin{tcolorbox}[breakable,colback=teal!4,colframe=teal!35!black,title={Before you move on}]
What does \bn{ওঠানামা} \bn{করা} mean? (\textquotedblleft{}To go up and down; to fluctuate\textquotedblright{}.) Say it once more.
\end{tcolorbox}
